\documentclass{article}
\usepackage{amsmath}
\usepackage{amsfonts}
\usepackage{amssymb}
\usepackage{bm}
\usepackage{hyperref}

\title{Exercise 2}
\author{}
\date{}

\begin{document}

\maketitle

\section*{Convex Functions (1 pt)}

\subsection*{Simple functions (0.5 pt)}

Are the following functions convex? Prove your statement.

\begin{enumerate}
    \item $f(x) = x^2$, $x \in \mathbb{R}$
    \item $f(x) = e^{x^2}$, $x \in \mathbb{R}$
    \item $f(x, y) = x^2 + 3xy + 2y^2$, $x \in \mathbb{R}$, $y \in \mathbb{R}$
\end{enumerate}

1.
\begin{align*}
             & f(x) = x^2                                          \\
    \implies & f'(x) = 2x                                          \\
    \implies & f''(x) = 2 > 0 \quad \forall x \in \mathbb{R} \quad \\
             & \therefore \text{$f$ is convex}
\end{align*}

2.
\begin{align*}
             & f(x) = e^{x^2}                                                                                   \\
    \implies & f'(x) = 2x e^{x^2}                                                                               \\
    \implies & f''(x) = 2 e^{x^2} + (2x) e^{x^2} (2x) = 2 e^{x^2} (2x^2 + 1) > 0 \quad \forall x \in \mathbb{R} \\
             & \therefore \text{$f$ is convex}
\end{align*}

3.
\begin{align*}
             & f(x, y) = x^2 + 3xy + 2y^2          \\
    \implies & \nabla f(x, y) = (2x + 3y, 3x + 4y) \\
    \implies & \nabla^2 f(x, y) =
    \begin{bmatrix}
        2 & 3 \\
        3 & 4
    \end{bmatrix}
\end{align*}

$\nabla^2 f(x, y)$ is not positive semi-definite since its eigenvalues are $3 + \sqrt{10}$ and $3 - \sqrt{10} < 0 $.
Therefore $f$ is not convex.

\subsection*{Log-sum-exp (0.5 pt)}

Show that the function
\[
    f(\bm{x}) = \log(e^{x_1} + \ldots + e^{x_n})
\]
is convex on $\mathbb{R}^n$.

\textit{Hint:} The proof is outlined in Boyd's book p. 74. Develop it yourself and provide the intermediate steps.

\textbf{Solution:} We will develop the proof in Boyd's book.

Let's define for convenience two column vectors:

\[
    z =
    \begin{bmatrix}
        e^{x_1} &
        e^{x_2} &
        \dots   &
        e^{x_n}
    \end{bmatrix}^{T}
\]

\[
    \bm{1} = \begin{bmatrix}
        1     &
        1     &
        \dots &
        1
    \end{bmatrix}^{T}
\]

This allows us to rewrite $f(\bm{x})$ in a more compact manner since:

\begin{align}
     & \bm{1}^{T} z = \sum_{i=1}^{n} e^{x_i}                        \\
     & f(\bm{x}) = \log(\sum_{i=1}^{n} e^{x_i}) = log(\bm{1}^{T} z)
\end{align}

First calculate the gradient:

\begin{align}
             & \frac{\partial f}{\partial x_i} =  \frac{e^{x_i}}{\sum_{i=1}^{n} e^{x_i}} \\
    \implies & \nabla f(\bm{x}) = \frac{z}{\bm{1}^{T} z}
\end{align}

Then calculate the Hessian:

%  
% 

\begin{align}
    \frac{\partial f}{\partial x_i^2} & = \frac{(-1) e^{x_i}}{\left( \sum_{i=1}^{n} e^{x_i} \right)^2} + \frac{e^{x_i}}{\sum_{i=1}^{n} e^{x_i}}                      \\
                                      & = \frac{1}{\left( \sum_{i=1}^{n} e^{x_i} \right)^2} \left( \left(\sum_{i=1}^{n} e^{x_i}\right) e^{x_i} - (e^{x_i})^2 \right)
\end{align}

\begin{align}
    \frac{\partial f}{\partial x_i \partial x_j} & = \frac{(-1) e^{x_i} e^{x_j}}{\left( \sum_{i=1}^{n} e^{x_i} \right)^2}
\end{align}

Notice that the outer product $ z z^T $ gives us the matrix:

\begin{equation}
    z z^T =
    \begin{bmatrix}
        e^{x_1} \\
        e^{x_2} \\
        \dots   \\
        e^{x_n}
    \end{bmatrix} \begin{bmatrix}
        e^{x_1} &
        e^{x_2} &
        \dots   &
        e^{x_n}
    \end{bmatrix} = \begin{bmatrix}
        (e^{x_1})^2     & e^{x_1} e^{x_2} & e^{x_1} e^{x_3} & \dots & e^{x_1} e^{x_n} \\
        e^{x_1} e^{x_2} & (e^{x_2})^2     & e^{x_2} e^{x_3} & \dots & \vdots          \\
        e^{x_1} e^{x_3} & e^{x_2} e^{x_3} & (e^{x_3})^2     & \dots & \vdots          \\
        \vdots          & \dots           & \dots           & \dots & \vdots          \\
        e^{x_1} e^{x_n} & \dots           & \dots           & \dots & (e^{x_n})^2
    \end{bmatrix}
\end{equation}

This accounts for the terms $- (e^{x_i})^2$ in $\frac{\partial f}{\partial x_i^2}$ and $(-1) e^{x_i} e^{x_j}$ in $\frac{\partial f}{\partial x_i \partial x_j}$.

Next, define a matrix with $z$ in its diagonal:

\begin{equation}
    \bm{diag}(z) = \begin{bmatrix}
        e^{x_1} & 0       & 0       & \dots & 0       \\
        0       & e^{x_2} & 0       & \dots & \vdots  \\
        0       & 0       & e^{x_3} & \dots & \vdots  \\
        \vdots  & \dots   & \dots   & \dots & \vdots  \\
        0       & \dots   & \dots   & \dots & e^{x_n}
    \end{bmatrix}
\end{equation}

This will take care of the terms $e^{x_i}$ in $\frac{\partial f}{\partial x_i^2}$.

Combining the previous results, we can express the gradient with compact notation as:

\begin{equation}
    \nabla^2 f(\bm{x}) = \frac{1}{(\bm{1}^{T} z)^2} \left( (\bm{1}^{T} z) \bm{diag}(z) - z z^T \right)
\end{equation}

To verify that $ \nabla^2 f(\bm{x}) \succeq 0 $ we must show that $ \forall v, v^T \nabla^2 f(\bm{x}) v \geq 0 $.
Distributing the $v$ and $v^T$ to the inside, we get:

\begin{align}
    v^T \nabla^2 f(\bm{x}) v & = \frac{1}{(\bm{1}^{T} z)^2} \left( (\bm{1}^{T} z) v^T \bm{diag}(z) v - v^T z z^T v \right)
\end{align}

We care only about the two terms inside the parentheses:

\begin{align}
    v^T z z^T v & = (v^T z) (z^T v)                                                                     \\
                & =
    \left(
    \begin{bmatrix}
            v_1   &
            v_2   &
            \dots &
            v_n
        \end{bmatrix}
    \begin{bmatrix}
            e^{x_1} \\
            e^{x_2} \\
            \dots   \\
            e^{x_n}
        \end{bmatrix}
    \right)
    \left(
    \begin{bmatrix}
            e^{x_1} &
            e^{x_2} &
            \dots   &
            e^{x_n}
        \end{bmatrix}
    \begin{bmatrix}
            v_1   \\
            v_2   \\
            \dots \\
            v_n
        \end{bmatrix}
    \right)                                                                                             \\
                & = \left( \sum_{i=1}^{n} v_i e^{x_i} \right) \left( \sum_{i=1}^{n} v_i e^{x_i} \right) \\
                & = \left( \sum_{i=1}^{n} v_i e^{x_i} \right)^2
\end{align}

And for the other term:

\begin{equation}
    v^T \bm{diag}(z) v = \sum_{i=1}^{n} v_i^2 e^{x_i}
\end{equation}

Replacing these results we get:

\begin{align}
    v^T \nabla^2 f(\bm{x}) v & = \frac{1}{(\bm{1}^{T} z)^2} \left( \left( \sum_{i=1}^{n} e^{x_i} \right) \left( \sum_{i=1}^{n} v_i^2 e^{x_i} \right) - \left( \sum_{i=1}^{n} v_i e^{x_i} \right)^2 \right) \geq 0
\end{align}

Take the Cauchy-Schwarz inequality:

\begin{equation}
    (\bm{a^T} \bm{a}) (\bm{b^T} \bm{b}) \geq (\bm{a^T} \bm{b})^2 \quad \forall \bm{a}, \bm{b} \in \mathbb{R}^n
\end{equation}

If $ a_i = v_i \sqrt{e^{x_i}} $ and $ b_i = \sqrt{e^{x_i}} $, we obtain:

\begin{align}
    \left( \sum_{i=1}^{n} v_i^2 e^{x_i} \right) \left( \sum_{i=1}^{n} e^{x_i} \right)                                               & \geq \left( \sum_{i=1}^{n} v_i e^{x_i} \right)^2 \\
    \left( \sum_{i=1}^{n} v_i^2 e^{x_i} \right) \left( \sum_{i=1}^{n} e^{x_i} \right) - \left( \sum_{i=1}^{n} v_i e^{x_i} \right)^2 & \geq 0
\end{align}

This last equation is the inside of $ v^T \nabla^2 f(\bm{x}) v $. Therefore $f$ is convex.

\subsection*{Geometric Mean (Bonus (0.5 pt))}

Show that the geometric mean
\[
    f(\bm{x}) = \left( \prod_{i=1}^{n} x_i \right)^{1/n}
\]
is concave on $\mathbb{R}_{++}^n$.

\textbf{Solution:} One possible way is to analyze the Hessian (Boyd's book p. 74) but there is \href{https://math.stackexchange.com/questions/3839629/how-to-prove-the-geometric-mean-is-concave/3839952#3839952}{a more elegant solution}.

We will use the AM-GM inequality: The arithmetic mean is greater or equal than the geometric mean.

\[
    \frac{1}{n} \sum_{i=1}^{n} a_i \geq \left( \prod_{i=1}^{n} a_i \right)^{1/n}
\]

Replace in the previous equation the following two values:

\[
    a_i = \frac{x_i}{\lambda x_i + (1 - \lambda) y_i} \quad b_i = \frac{y_i}{\lambda x_i + (1 - \lambda) y_i}
\]

To obtain two inequalities:

\begin{align}
     & \left( \prod_{i=1}^{n} \frac{x_i}{\lambda x_i + (1 - \lambda) y_i} \right)^{1/n} = \frac{f(x_i)}{f(\lambda x_i + (1 - \lambda) y_i)} \leq \frac{1}{n} \sum_{i=1}^{n} \frac{x_i}{\lambda x_i + (1 - \lambda) y_i} \label{am-gm-1} \\
     & \left( \prod_{i=1}^{n} \frac{y_i}{\lambda x_i + (1 - \lambda) y_i} \right)^{1/n} = \frac{f(y_i)}{f(\lambda x_i + (1 - \lambda) y_i)} \leq \frac{1}{n} \sum_{i=1}^{n} \frac{y_i}{\lambda x_i + (1 - \lambda) y_i} \label{am-gm-2}
\end{align}

Now multiply \eqref{am-gm-1} by $\lambda$ and \eqref{am-gm-2} by $(1 - \lambda)$:

\begin{align}
     & \frac{\lambda f(x_i)}{f(\lambda x_i + (1 - \lambda) y_i)} \leq \frac{1}{n} \sum_{i=1}^{n} \frac{\lambda x_i}{\lambda x_i + (1 - \lambda) y_i} \label{am-gm-3}             \\
     & \frac{(1 - \lambda) f(y_i)}{f(\lambda x_i + (1 - \lambda) y_i)} \leq \frac{1}{n} \sum_{i=1}^{n} \frac{(1 - \lambda) y_i}{\lambda x_i + (1 - \lambda) y_i} \label{am-gm-4}
\end{align}

Now sum both \eqref{am-gm-3} and \eqref{am-gm-4}:

\begin{align*}
             & \frac{\lambda f(x_i) + (1 - \lambda) f(y_i)}{f(\lambda x_i + (1 - \lambda) y_i)} \leq \frac{1}{n} \sum_{i=1}^{n} \left( \frac{\lambda x_i}{\lambda x_i + (1 - \lambda) y_i} + \frac{(1 - \lambda) y_i}{\lambda x_i + (1 - \lambda) y_i} \right) = 1 \\
             & \frac{\lambda f(x_i) + (1 - \lambda) f(y_i)}{f(\lambda x_i + (1 - \lambda) y_i)} \leq 1                                                                                                                                                             \\
    \implies & \lambda f(x_i) + (1 - \lambda) f(y_i) \leq f(\lambda x_i + (1 - \lambda) y_i)
\end{align*}

Now take $g(\bm{x}) = -f(\bm{x})$ and replace:

\[
    g(\lambda x_i + (1 - \lambda) y_i) \leq \lambda  g(x_i) + (1 - \lambda) g(y_i)
\]

Which is the definition of a convex function. Since $g(\bm{x}) = -f(\bm{x})$, then $f(\bm{x}) $ is a concave function.

\end{document}